\section{Conditions on a mechanism}\label{sec:conditions}

A candidate arising from this picture inherits the necessary conditions the
companion investigation has accumulated \cite{Selection}. We do not restate their
derivations; the table records which bear on a flow, and what this paper has
added to each.

\begin{center}
\begin{tabular}{@{}p{0.28\textwidth}p{0.64\textwidth}@{}}
\toprule
Inherited condition & Bearing, after this paper\\
\midrule
Reflection invariance & The flow must commute with $x\mapsto-x$; the involution is
the functional equation, which is also what makes
Proposition~\ref{prop:unimodular} true.\\
Unboundedness relative to $L^2$ & Now a statement about \emph{time}: by
Proposition~\ref{prop:groupdelay} the delay grows like $\log\tau$, so a packet at
frequency $N$ is held for $\sim\log N$. Condition~\ref{cond:bilinear} restates the
exclusion in the form a candidate can fail.\\
Spectral weight and channel tower & The archimedean multiplier \emph{is} the group
delay (Proposition~\ref{prop:groupdelay}), so the density identity is no longer an
input to be matched but a property of a phase.\\
Criticality & Section~\ref{sec:region}: the flow formulation inherits it. A bound
at fixed $\omega$ buys a zero-free strip of width $\omega$ and no more.\\
Jump structure & Realized on the connection by \eqref{eq:atomic}: the atoms of the
kernel are the prime periods.\\
\bottomrule
\end{tabular}
\end{center}

Three conditions are specific to the deformation picture. They are
\emph{proposals}: each is elementary, none has been tested against a specific
theory, and the third carries an identification hypothesis this program
deliberately does not make.

\subsection{The object is a bilinear, not a character}

\begin{condition}[Bilinear endpoints, and a delay test]\label{cond:bilinear}
The object must be a bilinear in unbounded endpoint operators joined by a
transport --- not a character --- and its group delay must grow like $\log\tau$
with smooth part $2\theta'(\tau)$.
\end{condition}

The first clause names a class and the second is the test. The exclusion behind
the clause is worth restating precisely, because it is narrower than it looks. A
holonomy valued in a compact group, in a finite-dimensional representation, is
unitary, so its character is a bounded function and every pairing assembled from
it is bounded on $L^2(I_L)$ --- excluded by the unboundedness rule. That argument
uses only that the object is a \emph{trace of a group element}. It is also the
structural reason the conformal Wilson benchmark of \cite{Selection} failed:
$2\cos\theta$ is bounded \emph{because} it is a character of $SU(2)$ in the
fundamental, not by accident of that benchmark.

An \emph{open} Wilson line is not a character. In the construction of
\cite{OWL} the operator is
\[
 OW_i^{\,j}[\tilde C]=\bar\psi^{\,j}(x_j)\;
 P\exp\!\left[\int_{\tilde C}\!ds\,\Big(iA_\mu\dot x^\mu
 +n^k(s)\Phi_k\,|\dot x|\Big)\right]\psi_i(x_i),
\]
a bilinear in matter fields joined by a Wilson line, the endpoint fields
transforming in the fundamental and antifundamental so that the contraction is
gauge invariant. Its endpoints carry operator-valued distributions, so the object
is unbounded by construction and the character argument has nothing to act on.
\textbf{The condition is satisfied structurally rather than evaded}, which is why
it is stated here as a prescription rather than a prohibition.

Two further remarks connect that construction to Section~\ref{sec:endpoint}.
First, the geometry matches: the supersymmetric configuration of \cite{OWL} in
the defect theory is a \emph{ray ending on a defect}, one end pinned and one end
free, and Proposition~\ref{prop:endpoint} found the variation to be rank one at
the leading end only, the trailing one killed by causality --- with the boundary
value of the evolved field as the coefficient, which in the open Wilson line is
the endpoint field itself. Second, the obstruction of Proposition~\ref{prop:trace}
is, in that language, the ordinary fact that a field at a point is a distribution
and must be smeared; Condition~\ref{cond:marginal} fixes the profile.

Note also that $V_{\omega,L}$ is not unitary --- it is a contraction and the
object of interest is its defect \eqref{eq:defect} --- so the transport to be
matched is not anti-Hermitian either.

\subsection{The non-commutativity must be made visible}

\begin{condition}[Visible transverse structure]\label{cond:visible}
A mechanism must supply a reason for the holonomy to see the non-commutativity of
Proposition~\ref{prop:algebra}: either a domain that is not simply connected, or a
failure of the connection to exist.
\end{condition}

This replaces the weaker demand of Remark~\ref{rem:cost}, which asked only for
non-commutativity transverse to $\omega$. That exists
(Proposition~\ref{prop:algebra}); what it lacks is visibility
(Proposition~\ref{prop:flat}). Problems~\ref{prob:trace} and \ref{prob:curvature}
are the two forms the condition can take.

\subsection{Marginal short distance: effective dimension one half}

\begin{condition}[Marginality]\label{cond:marginal}
Suppose the line's parameter is identified with the arithmetic coordinate $x$.
Since $\gammak(r)=\frac1{2r}+O(1)$ as $r\downarrow0$ by \eqref{eq:gamma-kernel},
\eqref{eq:gamma-energy} gives
\[
 K[F]=\frac12\iint|F(x)-F(y)|^2\gammak(|x-y|)\dd x\dd y
 \sim\frac14\iint\frac{|F(x)-F(y)|^2}{|x-y|}\dd x\dd y
\]
near the diagonal: the Gagliardo seminorm at the borderline exponent $s=0$. An
insertion on the line reproducing this must have two-point function $|x-y|^{-1}$,
an effective defect dimension $\Delta=\frac12$, with the coefficient fixed.
\end{condition}

This excludes the first thing one would try. The displacement operator of a line
defect has protected dimension $\Delta=2$, fixed by the Ward identity that defines
it, and the standard protected scalar insertions on a supersymmetric Wilson line
have $\Delta=1$ \cite{GiombiKomatsu,CHMS}. \emph{So the deformation whose
generator matches is not the displacement operator}, which is what a deformation
of the contour would naively produce.

\begin{remark}[The principal series, and why everything here is marginal]
\label{rem:principal}
There is one reading under which Condition~\ref{cond:marginal} is not an accident.
For a one-dimensional conformal system the dilatation eigenvalue is the conformal
weight, and the unitary principal series of $SL(2,\R)$ is $\Delta=\frac12+i\tau$
with $\tau$ real --- the critical line. In that reading the zeros are
principal-series data, the shift $\omega$ moves off the unitary line, which is the
detection threshold of Corollary~\ref{cor:graded}, and the Riemann hypothesis is
the statement that nothing leaves it. Condition~\ref{cond:marginal}'s
$\Delta=\frac12$ is then the \emph{edge} of the principal series, the most
marginal unitary weight there is, which would explain the coincidence of
Remark~\ref{rem:marginal}: the Gagliardo borderline, the transfer's
$(2+|\tau|)^{-1/2}$ gain and the trace threshold $r>\frac12$ are three faces of
one exponent. Section~\ref{subsec:conformal} supplies the missing motivation: on a ray, the
Mellin variable is the dilation eigenvalue and $s=\Delta+i\tau$, so
$\Delta=\frac12$ puts it on the critical line --- and $\frac12$ is what
Condition~\ref{cond:marginal} requires for an unrelated reason. The identification
of the arithmetic coordinate with $\log r$ is therefore \emph{derived from a
symmetry match} rather than assumed: if one asks for a geometry whose residual
conformal symmetry is the Weil form's exact symmetry group, the ray with a defect
at one end is that geometry. \emph{The program's caveat is not removed --- this
remains an identification --- but there is now a reason to make it. The
neighbourhood, conformal quantum mechanics and the zeros, is crowded.}
\end{remark}

\subsection{Conformal symmetry: the form already has it}\label{subsec:conformal}

It is natural to ask whether the conformal symmetry of Yang--Mills could be used
to compactify and tame the infinities of the problem. The answer is better than
the question expects, and it corrects a claim made in version 0.2 of this
manuscript.

\paragraph{The form's exact symmetry group is a conformal one.}
Every term of \eqref{eq:target} is either a multiple of $\lVert f\rVert^2$ or an
integral against a kernel depending on $x-y$, so the Weil functional is invariant
under $f\mapsto f(\cdot-a)$, exactly and unconditionally --- the Toeplitz
structure the companion investigation's numerics exploit. Write $x=\log r$, so
that the arithmetic coordinate is the logarithm of a radial coordinate on a ray.

\begin{center}
\begin{tabular}{@{}lll@{}}
\toprule
in $x$ & in $r$ & status for the Weil form\\
\midrule
translation $x\mapsto x+a$ & dilation $r\mapsto\lambda r$ & exact symmetry\\
reflection $x\mapsto-x$ & inversion $r\mapsto1/r$ & exact symmetry (rule 1)\\
dilation $x\mapsto\lambda x$ & $r\mapsto r^\lambda$ & not a symmetry, and not conformal\\
\bottomrule
\end{tabular}
\end{center}

The Möbius maps of the line preserving $\{0,\infty\}$ are exactly
$r\mapsto\lambda r$ and $r\mapsto\mu/r$, so the stabilizer of a ray's two
endpoints is generated by dilation and inversion --- and the Weil form is
invariant under all of it. The special conformal generators move the endpoints
and are broken by the configuration, not by the form. Hence:

\begin{center}
\emph{the Weil form's exact symmetry group is precisely the residual conformal
group of a ray with one end pinned at the origin and the other at infinity,}
\end{center}

which is the configuration of \cite{OWL}. Both halves of that symmetry were
already program facts; what is new is that they are the stabilizer of a geometry.

\paragraph{Correction.} Version 0.2 argued here that exact scale invariance is
incompatible with the target, citing the exclusion of \cite{ExclusionMap} that
dilation destroys the prime atoms. That exclusion concerns dilation \emph{in
$x$}, which is $r\mapsto r^\lambda$ and is not a Möbius map. Under the conformal
dilation $r\mapsto\lambda r$ the atoms sit at fixed \emph{ratios}
$r_2/r_1=p^m$ and are preserved. In the ray picture the Weil form has no absolute
scale to protect, and conformal invariance is not in conflict with it. What the
atoms require is a preferred set of ratios, which is a statement about the
spectrum of dimensions rather than about scale: in Mellin variables the prime
part carries weight $-\zeta'/\zeta$, a continuum, and a conformal system with a
continuous spectrum of dimensions is exactly what the principal series is. The
demand is therefore not \emph{break conformal invariance} but \emph{carry this
continuum of dimensions}, which is sharper and more useful. We do not challenge
the exclusion of \cite{ExclusionMap} in its own setting, where the relevant
dilation is the field theory's own.

\paragraph{Positivity is a conformal-orbit invariant.}
\begin{proposition}[Congruence]\label{prop:congruence}
If $M$ is invertible on the form domain then $Q\succeq0$ if and only if
$M^*QM\succeq0$. A conformal reparametrization acts on a weight-$\Delta$ kernel
by congruence with multiplication by a positive power of the Jacobian, so
positivity of the target is invariant along its conformal orbit.
\end{proposition}
This is Sylvester's law of inertia. Since the criterion asks only for a sign, any
member of the conformal orbit of $Q_L$ is an equally good target, and the
ray-to-semicircle map of \cite{OWL} is not merely a reframing: the semicircle
form is a genuinely different quadratic form with the same positivity content, on
a \emph{compact} contour. The program already uses one congruence of this kind,
the compression presentation $Q_L=K_+-T_L$ read as $T_L\preceq K_+$; conformal
congruences are a second, geometrically natural class that has not been used.

Two guards. The freedom is vacuous if the class is unrestricted, since every
positive form is congruent to the identity --- any use must name the class in
advance. And the Jacobian of the ray-to-semicircle map \emph{degenerates at the
endpoints}, so the congruence is boundedly invertible in the interior and fails
exactly where Proposition~\ref{prop:trace} already obstructs. Whether that
degeneration softens the endpoint or leaves the form domain is open, and it is the
first thing to check if this route is pursued.

\paragraph{What compactifies.} The spectral infinity does, and the
compactification is already in hand: the conformal map of the right half-plane to
the disc is the Cayley transform of Proposition~\ref{prop:cayley}, under which
$K_\omega$ becomes an inner function of the disc whose zeros accumulate at one
boundary point, the logarithmically divergent delay of
Proposition~\ref{prop:groupdelay} being exactly that accumulation. The length
infinity does too, by the ray-to-semicircle map. What does not is
Proposition~\ref{prop:trace}, a short-distance statement invariant under
reparametrization; and Corollary~\ref{cor:graded}, a statement about the zero set
of $\xi$, so compactifying moves its content to a boundary point without removing
it.

\paragraph{The vanishing theorem, read the other way.}
\cite{OWL} finds that in non-superconformal $\mathcal N=2$ theories with
fundamental hypermultiplets all BPS open Wilson lines have vanishing expectation
value, nontrivial protected operators appearing only in the superconformal case.
Version 0.2 read this as a trap. It is not: superconformality is available and is
not in conflict with the target, because the target's exact symmetry group is the
residual conformal group of the geometry in which those operators live. That
makes the superconformal defect setting the natural place to run the test of
Condition~\ref{cond:bilinear}, rather than a place where it was expected to fail
for structural reasons.

\paragraph{Where the ratios could come from.} Not from a single conformal defect,
which preserves dilations about a point on it and therefore fixes no ratio.
Ratios are \emph{separations}: the atoms sit at $r_2/r_1=p^m$, so what is wanted
is a spectrum of dimensions carrying weight $-\zeta'/\zeta$, equivalently a
structure invariant under $r\mapsto\lambda r$ but distinguishing the ratios
$p^m$. That is \eqref{eq:atomic} in geometric dress, with the punctures at
integer ratios, and it is the same statement as Condition~\ref{cond:visible}. The
named object in that direction is the Bost--Connes system \cite{BostConnes}; we
record where the structure lives and do not enter that literature.
